\section{Conclusion}
\label{sec:conclusion}

\paragraph{Objectifs.} Ce projet visait à reconnaître automatiquement quatre maladies foliaires de la canne à sucre, ainsi que les feuilles saines, à partir de photographies prises au champ, et à mesurer ce qu'apporte l'apprentissage par transfert lorsque les données sont peu nombreuses : quelle architecture choisir, jusqu'où adapter un réseau pré-entraîné, et s'il est utile de combiner plusieurs réseaux.

\paragraph{Méthodologie.} Après avoir audité le jeu de données de Daphal et Koli et retiré 234 images redondantes, nous avons comparé quatre architectures récentes et de tailles différentes, EfficientNetV2-S, MobileNetV4-Conv-M, ConvNeXtV2-T et SE-ResNet-50, chacune selon trois stratégies (extraction de caractéristiques, fine-tuning partiel et fine-tuning complet), ainsi qu'un CNN entraîné from scratch. Tous les modèles ont été entraînés avec les mêmes hyperparamètres, les mêmes lots et les mêmes augmentations, puis évalués par validation croisée à cinq plis sur un ensemble de test isolé, avec un bootstrap pour juger de la significativité des écarts. Nous avons également étudié la fusion des embeddings des quatre réseaux sur toutes les combinaisons possibles, le vote souple, la robustesse aux dégradations d'image, le coût de calcul, le besoin en données et l'interprétation des décisions par Grad-CAM.

\paragraph{Meilleurs résultats.} Le fine-tuning complet de ConvNeXtV2-T donne le meilleur modèle unique, avec $98{,}9 \pm 0{,}6$\,\% d'accuracy, un F1 macro de 98,9\,\% et une AUC de 0,9998 ; c'est aussi le plus robuste aux dégradations et celui dont les cartes Grad-CAM suivent le mieux les lésions. La fusion des quatre réseaux fine-tunés, par concaténation des embeddings ou par vote souple, atteint 99,3 à 99,4\,\% et réduit de moitié environ le taux d'erreur, au prix d'un calcul trois fois plus lourd. Pour un usage sur téléphone, MobileNetV4-M offre le meilleur compromis : 97,9\,\% d'accuracy avec 8,4 millions de paramètres et 30~ms par image sur CPU, à condition de surveiller sa sensibilité au bruit.

\paragraph{Apport du transfer learning.} L'apprentissage par transfert est nettement gagnant, mais à une condition : il faut fine-tuner le réseau. Utilisés comme simples extracteurs de caractéristiques, les réseaux ImageNet plafonnent entre 84 et 91\,\%, en dessous d'un petit CNN entraîné from scratch (94,3\,\%). Une fois fine-tunés, ils divisent par cinq le nombre d'erreurs de ce CNN, convergent en quelques époques, résistent mieux aux dégradations d'image et concentrent leur attention sur les symptômes. Ils ont aussi besoin de beaucoup moins de données : avec un quart des images d'apprentissage, ConvNeXtV2 dépasse déjà le CNN from scratch entraîné sur la totalité, et avec 155 images seulement il conserve 91,5\,\% d'accuracy, contre 71,5\,\% pour ce dernier. Le pré-entraînement sur ImageNet apporte donc surtout un excellent point de départ, que le fine-tuning adapte ensuite aux textures et aux couleurs des lésions foliaires.

Ces résultats ont été obtenus sur un jeu de données acquis dans une seule région et dans des conditions assez homogènes. La prochaine étape consisterait à les confirmer sur des photographies prises ailleurs, avec d'autres téléphones et à d'autres saisons, avant d'envisager un déploiement de MobileNetV4 sur smartphone, puis le passage à la détection des maladies dans des images de drone.
